\section{Authority: is this within Congress's power and consistent with existing law?}
\label{sec:authority}

The second question is whether Congress may act here at all, and whether the bill
respects the law already on the books. A clean answer disarms both the
constitutional objection and the federalism objection.

[DRAFTING INSTRUCTIONS: In the full stage, ground authority in the bill's structure
from \texttt{s3-amendment.tex} and \texttt{s4-comparative.tex}: the bill amends the
Federal Food, Drug, and Cosmetic Act (21 U.S.C. 301 et seq.), an exercise of the
commerce power over medical devices; it adds section 515D and ten conforming
amendments; its rule of construction leaves the investigational-device exemption
(520(g)), human-subjects protection (21 CFR part 50), and the IND framework (21 CFR
part 312) undisturbed; and its savings clause (new 360k(c)) preserves State
human-over-AI laws. Cite \cite{fdcact}, \cite{cures2016}, and the closest analogues
(section 515C predetermined change control plans; the CMS Medicare Advantage AI rule
\cite{cms2024aiguardrails}). Use the author's regulatory adaptations
\cite{Kawchak2026Adaption21CFRPart, Kawchak2026Adaption21CFRPartb,
Kawchak2026AdaptionICHHarmonisedGuideline}. This section carries the amendment-map
figure (diagram 05) and the federalism figure (diagram 09).]

\begin{psgfig}
\node[psgevid] (fdc) at (0,0) {FD\&C Act};
\node[psgact] (s) at (4.0,0) {Section 515D\\floor};
\node[psgthree] (save) at (8.0,0) {State savings\\clause};
\draw[psgedge] (fdc)--(s); \draw[psgedge] (s)--(save);
\end{psgfig}
\psgcaption{Figure (placeholder). Authority flows from the FD\&C Act to a Federal
floor that preserves State human-over-AI law. [FULL STAGE: replace with diagram 09.]}

[FULL STAGE: add Table (authority) mapping each rule-of-construction clause to the
existing-law provision it preserves, body width, L/Y columns.]
